\chapter{Create and submit the course repository}
\label{ch:repo-submission}

\subsection{Create an empty GitHub repository}

Choose one method below to create an empty remote repository: use the GitHub website steps or the optional terminal commands. Do not complete both methods.

\begin{enumerate}[leftmargin=1.7em,itemsep=0.35em,topsep=0.4em]
  \item Sign in to GitHub. In the upper-right corner, select the \textbf{+} menu and choose \textbf{New repository}.
  \item Choose the account or organization that should own the repository.
  \item Enter a repository name, such as \texttt{UTP-course-name}, and select the visibility required by your instructor.
  \item Leave \textbf{Add a README file} unchecked, and leave the \textbf{.gitignore} template and license set to \textbf{None}.
  \item Select \textbf{Create repository}. Keep the new repository page open; you will connect it to your local project in the next section.
\end{enumerate}

\subsection{Prepare the local project and connect it to GitHub}

Do not upload the virtual environment, passwords, tokens, or API keys. In the project root, create a file named \texttt{.gitignore}.

From the project root, create the first commit and connect it to the empty
GitHub repository. Replace \texttt{USERNAME} and \texttt{REPOSITORY}:

\begin{lstlisting}[style=utpterminal]
git init -b main
git add .
git status
git commit -m "initial commit"

git remote add origin git@github.com:USERNAME/REPOSITORY.git
git remote -v
git push -u origin main
\end{lstlisting}

\begin{pitfall}
Always read \texttt{git status} before committing. If a dataset, password,
token, \texttt{.env} file, or virtual environment appears, stop and update
\texttt{.gitignore} before running \texttt{git commit}.
\end{pitfall}